% \chapter{Pruning i budowa architektury MoE}

% \section{Próba lokalizacji i izolacji wąskiej dziedziny wiedzy}

% \section{Przeprowadzenie pruningu wag nieskorelowanych w celu uzyskania mniejszego, domenowego modelu}

% \section{Implementacja klasyfikatora sterującego (routera) bazującego na zamrożonych cechach}

% \section{Substytucja gęstych warstw zespołem wyspecjalizowanych ekspertów (hard-routing)}


\chapter{Pruning i budowa architektury MoE}
\label{chap:pruning_moe}
Mapowanie opisane w rozdziale~\ref{chap:klasteryzacja} dostarcza informacji o związku neuronów MLP z sygnałami domen wyodrębnionych przez SAE. Kolejnym krokiem jest wykorzystanie tych zależności do modyfikacji modelu bazowego. W niniejszym rozdziale przedstawiono tworzenie przyciętych (ang. \textit{pruned}) modeli domenowych, trening klasyfikatora sterującego oraz połączenie obu elementów w architekturę mieszaniny ekspertów. Rozwiązanie stanowi konwersję już wytrenowanego modelu, nawiązującą do podejścia \textit{MoEfication} \cite{zhang_2022_moefication}, omówionego w podrozdziale~\ref{sec:moefication}. Eksperci powstają przez wybór neuronów istniejącej warstwy, a uczonym elementem jest wyłącznie router. Szczegółowe porównanie jakości modeli po modyfikacji stanowi przedmiot rozdziału~\ref{chap:walidacja}.

\section{Zakres izolacji domen i konstrukcja eksperymentu}
\label{sec:izolacja_domen}
Próbę wydzielenia modułów dziedzinowych ograniczono do MLP warstwy $9$ modelu \textit{Gemma 3 270M}. Jest to ta sama warstwa, dla której wykonano analizę SAE i mapowanie aktywacji domenowych na neurony MLP. Pozostałe bloki transformera, mechanizm uwagi oraz reprezentacje wejściowe i wyjściowe modelu zachowują bazowe parametry. \newline
Umiejscowienie ekspertów w części MLP odpowiada rozwiązaniom stosowanym w architekturach mieszaniny ekspertów, opisanym w podrozdziale~\ref{sec:moe_transformer}. W rozbudowanych modelach mechanizm ten obejmuje jednak wiele warstw, a zależnie od architektury może występować we wszystkich blokach transformera. W niniejszej pracy zakres modyfikacji celowo ograniczono do jednej warstwy, aby w skali dostępnego eksperymentu ocenić skutki semantycznego pruningu i warunkowego wyboru ekspertów. Nie zakłada się przy tym, że pełna wiedza o danej dziedzinie znajduje się tylko w tej warstwie. Badanie ma ustalić, czy wyodrębnione zależności między domenami SAE a neuronami MLP mogą posłużyć do ograniczenia aktywnej szerokości modułu oraz ilości wykonanych obliczeń przy zachowaniu jakości modelu. \newline
Dla każdej z sześciu domen wybranych w podrozdziale~\ref{sec:wybor_domen} tworzony jest osobny wektor maski o długości równej szerokości pośredniej MLP. Maska określa, które neurony zostają zachowane w danym ekspercie. Nie jest przeprowadzany także dodatkowy treningu wag ekspertów ani dostrajanie modelu językowego po pruningu. Określenie \enquote{ekspert domenowy} oznacza zatem podmoduł wybrany na podstawie związku aktywacji z domeną, a nie niezależnie wytrenowany model specjalistyczny. \newline
Zaimplementowane zostały dwa sposoby wykorzystania takiego podmodułu. W wariancie pojedynczego eksperta przycięty MLP zastępuje bazowy moduł dla wszystkich tokenów. W wariancie MoE decyzja jest podejmowana osobno dla każdej pozycji: router wybiera jednego z ekspertów albo kieruje token do oryginalnego MLP. Pierwszy wariant pozwala ocenić skutki samej selekcji neuronów, drugi zaś sprawdza możliwość warunkowego stosowania tej selekcji. \newline
Dobór progów, szerokości ekspertów i konfiguracji routera odbywa się na zbiorze uczącym. Końcowy zbiór testowy pozostaje oddzielony od tych decyzji. Aby odróżnić użyteczność kryterium semantycznego od skutków samej redukcji szerokości, w procedurze walidacyjnej przygotowane zostały również maski losowe i maski oparte na wielkości wag, zachowujące tę samą liczbę neuronów.

\section{Pruning neuronów na podstawie korelacji domenowych}
\label{sec:pruning_korelacyjny}

Proces pruningu korzysta z korelacji $r_{j,d}$ wyznaczonych w poprzednim etapie (podrozdział~\ref{sec:mapowanie_domen_mlp}). Wynikiem oceny neuronu $j$ dla domeny $d$ jest $|r_{j,d}|$, a więc siła zależności liniowej niezależnie od jej znaku. Próg zachowania neuronu wyznacza się osobno dla każdej domeny, ponieważ rozkłady sygnałów klastrów i ich związki z MLP mogą się różnić. \newline
W zastosowanym wariancie próg oparto na empirycznym rozkładzie zerowym uzyskanym przez permutacje sygnału SAE. Wspólną macierz aktywacji MLP pozostawia się bez zmian, natomiast sygnał domenowy dzieli na kolejne bloki po $32$ tokeny i losuje kolejność bloków. Porządek wartości wewnątrz bloku zostaje zachowany, co częściowo uwzględnia lokalną zależność kolejnych tokenów. Dla każdej ze $100$ permutacji oblicza się korelacje ze wszystkimi neuronami i zapisuje największą wartość bezwzględną:
\[
M_d^{(b)}=\max_{1\leq j\leq H}
\left|\operatorname{corr}\left(a_j,\widetilde{s}_d^{(b)}\right)\right|,
\qquad
\tau_d=Q_{0{,}95}\left(M_d^{(1)},\ldots,M_d^{(100)}\right),
\]
gdzie $H=2048$, $\widetilde{s}_d^{(b)}$ oznacza przestawiony sygnał domeny, a $Q_{0{,}95}$ -- empiryczny kwantyl rzędu $0{,}95$. Zastosowanie maksimum uwzględnia jednoczesne porównywanie wielu neuronów: wynik rzeczywisty musi przekroczyć poziom wyznaczony przez najsilniejsze przypadkowe zależności w permutowanych danych. \newline
Na tej podstawie powstaje wstępna maska:
\[
m_{j,d}^{\mathrm{stat}}=\mathbf{1}\{|r_{j,d}|\geq\tau_d\}.
\]
Statystyka maksimum służy ograniczeniu ryzyka przypadkowego wyboru neuronu w obrębie domeny, lecz jej interpretacja zależy od przyjętego schematu permutacji. Bloki tokenów nie odtwarzają pełnej struktury zależności w dokumentach i mogą przecinać ich granice. Otrzymany próg jest więc kryterium empirycznym, a nie bezwarunkową gwarancją poziomu błędu statystycznego. Tym bardziej nie rozstrzyga on przyczynowej roli neuronu: niska korelacja nie dowodzi, że neuron jest zbędny dla poprawnego przetwarzania tekstu. \newline
Maski oparte wyłącznie na tym progu zachowywały od $1272$ do $1408$ neuronów. Próby na danych rozwojowych wykazały jednak nadmierny spadek jakości części ekspertów. Wprowadzono zatem dolne ograniczenie szerokości, sterowane parametrem \texttt{min\_keep\_fraction}. Jeżeli maska statystyczna zachowuje mniej niż $\lceil\rho H\rceil$ neuronów, wybierane jest tyle neuronów o najwyższych skończonych wartościach $|r_{j,d}|$. W końcowej konfiguracji przyjęto $\rho=0{,}75$, co dla każdej domeny dało $1536$ zachowanych i $512$ usuniętych neuronów. Wszystkie maski statystyczne były mniejsze od tego limitu, dlatego końcowy wybór był w praktyce wyborem najlepszych $1536$ neuronów według rankingu korelacyjnego. Nie należy zatem interpretować wszystkich zachowanych neuronów jako przekraczających próg permutacyjny. \newline
Maskowanie musi obejmować całą oś pośrednią MLP. W Gemmie ten sam indeks neuronu wyznacza wiersz macierzy \texttt{gate\_proj}, wiersz macierzy \texttt{up\_proj} i kolumnę macierzy \texttt{down\_proj}. Dla zbioru zachowanych indeksów $I_d$ odpowiadający mu kompaktowy ekspert ma postać:
\[
E_d(u)=W_{\mathrm{down}}[:,I_d]
\left(\operatorname{act}\left(W_{\mathrm{gate}}[I_d,:]u\right)
\odot W_{\mathrm{up}}[I_d,:]u\right).
\]
Wymiar wejścia i wyjścia pozostaje równy $640$, natomiast wymiar pośredni maleje z $2048$ do $1536$. Dzięki temu ekspert może zastąpić bazowy MLP bez zmiany interfejsu bloku transformera. \newline
Należy przy tym rozróżnić zapis maski od fizycznego zmniejszenia macierzy. Moduł pruningu eksportuje maski boolowskie oraz pełnowymiarowe stany MLP z wyzerowanymi odpowiednimi wierszami i kolumnami. Samo zerowanie nie zmniejsza kosztu standardowych operacji liniowych. Dopiero moduł składania modelu wybiera zachowane wycinki wag i tworzy mniejsze projekcje. Redukcja o $25\%$ dotyczy zatem szerokości i wag projekcji pojedynczego MLP eksperta, a nie całkowitej liczby parametrów modelu językowego.

\section{Trening routera na zamrożonych reprezentacjach modelu}
\label{sec:trening_routera}

Router został zaimplementowany w module \texttt{router\_training.py} jako pojedyncza warstwa liniowa o wymiarach $640\rightarrow6$. Jego wejściem jest wektor $u(t)$ przekazywany do MLP warstwy $9$, przechwytywany przed wykonaniem tego modułu. Jest to reprezentacja dostępna przed wyborem eksperta, odmienna od strumienia rezydualnego po bloku, używanego wcześniej do kodowania SAE. Router wyznacza logity i rozkład po domenach:
\[
\ell(t)=W_{\mathrm R}u(t)+b_{\mathrm R},
\qquad
p_d(t)=\frac{\exp(\ell_d(t))}{\sum_{k=1}^{6}\exp(\ell_k(t))}.
\]
Parametry Gemmy są zamrożone, a gradienty aktualizują wyłącznie $W_{\mathrm R}$ i $b_{\mathrm R}$. SAE służył do odkrycia domen i przygotowania masek; nie jest uruchamiany podczas treningu routera ani późniejszej inferencji MoE. \newline
Reprezentacje wejściowe zbierane są jednokrotnie dla każdej partii tekstów, bez obliczania gradientów modelu bazowego, a następnie przechowywane na CPU w formacie FP16. Kolejne epoki routera korzystają z tego bufora, ograniczając koszt powtarzania obliczeń Gemmy. Po usunięciu paddingu każdemu tokenowi przypisywana jest etykieta domeny jego tekstu. Aby długie dokumenty nie dominowały nad krótkimi, token z tekstu $i$ o długości $n_i$ otrzymuje wagę $w_t=1/n_i$. Dla partii tokenów $B$ minimalizowana jest ważona entropia krzyżowa:
\[
\mathcal{L}_B=
-\frac{\sum_{t\in B}w_t\log p_{y(t)}(t)}{\sum_{t\in B}w_t},
\]
gdzie $y(t)$ oznacza etykietę tekstu. Łączna waga tokenów danego tekstu w zbiorze wynosi jeden. Taki sposób ważenia wyrównuje wpływ długości dokumentów, ale nie usuwa niejednoznaczności etykiet: token ogólny lub początkowy fragment sekwencji nie musi jeszcze dostarczać wystarczającej informacji o domenie. \newline
Podział danych przeprowadzono na poziomie grup źródłowych, tak aby teksty należące do tej samej grupy nie znalazły się równocześnie w części treningowej i walidacyjnej. Po podziale uzyskano $860$ tekstów treningowych i $220$ walidacyjnych. Nie wykorzystano do tego celu diagnostycznych kontekstów, na podstawie których odkrywano klastry SAE. Router optymalizowano algorytmem AdamW przez $15$ epok, z tempem uczenia $10^{-3}$ i współczynnikiem regularyzacji wagowej $10^{-4}$. Najlepszy stan wybierano według tokenowego macro-F1 na części walidacyjnej; w wykonanym przebiegu był to stan z ostatniej epoki. \newline
Router nie ma osobnej klasy odpowiadającej tekstowi ogólnemu. Wprowadzono zamiast niej powrót do bazowego MLP (ang. \textit{fallback}), gdy największa wartość softmax jest zbyt niska. Próg $\theta$ skalibrowano na osobnym zbiorze $200$ ogólnych tekstów rozwojowych, przyjmując za cel skierowanie do ekspertów około $5\%$ ich tokenów. Wykorzystano empiryczny kwantyl maksymalnych wartości softmax; uzyskany próg wyniósł w przybliżeniu $0{,}8247$. Jest on zapisywany razem z wagami routera i kolejnością domen. \newline
Maksimum softmax pełni tutaj funkcję wskaźnika pewności, a nie gwarantowanego prawdopodobieństwa poprawnej decyzji. Kalibracja ogranicza routing na konkretnym korpusie ogólnym, lecz nie zapewnia takiego samego udziału dla dowolnych nowych danych. Wysoka rozróżnialność domen może ponadto częściowo wynikać ze stylu i pochodzenia tekstów. Z tego względu zarówno jakość klasyfikacji, jak i skutki wyboru eksperta wymagają późniejszej oceny na oddzielonym zbiorze końcowym.

\section{Złożenie ekspertów i tokenowy wybór modułu MLP}
\label{sec:skladanie_moe}

Moduł \texttt{moe\_assembly.py} odtwarza kompaktowych ekspertów i zastępuje MLP warstwy $9$ modułem \texttt{HardRoutedMLP}. Każdy ekspert zachowuje oryginalną funkcję aktywacji i kolejność projekcji, różniąc się zbiorem zachowanych neuronów. Bazowy, gęsty MLP pozostaje dostępny jako wariant ogólny. Wszystkie parametry są podczas inferencji zamrożone; nie prowadzi się wspólnego treningu routera i ekspertów po złożeniu modelu. \newline
Dla każdego tokenu wybierana jest domena o najwyższej wartości softmax. Jeżeli pewność przekracza lub osiąga skalibrowany próg, wykonywany jest odpowiadający jej ekspert; w przeciwnym razie stosowany jest bazowy moduł $E_{\mathrm{baz}}$:
\[
d^*(t)=\operatorname*{arg\,max}_{d}p_d(t),
\qquad
\operatorname{MLP}_{\mathrm{MoE}}(u(t))=
\begin{cases}
E_{d^*(t)}(u(t)), & \max_d p_d(t)\geq\theta,\\
E_{\mathrm{baz}}(u(t)), & \max_d p_d(t)<\theta.
\end{cases}
\]
Jest to twardy wybór jednego eksperta (ang. \textit{hard top-1 routing}): wyjścia ekspertów nie są sumowane ani ważone prawdopodobieństwami routera. Decyzja jest tokenowa, dlatego różne fragmenty tego samego tekstu mogą zostać przetworzone przez różne moduły. \newline
W implementacji tensor wejściowy jest spłaszczany do macierzy tokenów, a predykcje routera wyznaczają grupy pozycji przypisanych do poszczególnych ekspertów i wariantu ogólnego. Każdy moduł wykonuje obliczenia tylko dla swojej grupy, po czym wyniki są umieszczane z powrotem na właściwych pozycjach. Nie oblicza się wyjść wszystkich ekspertów dla całej partii. Pozostałe elementy warstwy oraz dalsze bloki modelu przetwarzają wynik w zwykły sposób. Powrót do bazowego MLP przywraca oryginalną funkcję tego modułu dla otrzymanego wejścia, ale nie gwarantuje identyczności całej predykcji z modelem bazowym, ponieważ reprezentacje mogą zależeć od wcześniejszych tokenów przetworzonych przez ekspertów. \newline
Maski domenowe mogą się nakładać. Wspólne neurony są kopiowane do odpowiednich ekspertów, a ich zachowane wagi nie są dodatkowo uczone. Architektura nie stanowi więc rozłącznego podziału parametrów na sześć części. Choć pojedynczy ekspert jest mniejszy od bazowego MLP, cały bank sześciu ekspertów wraz z pełnym wariantem ogólnym zwiększa pamięć potrzebną na parametry. Przy jednakowej szerokości ekspertów równej $0{,}75H$ ich projekcje wraz z projekcjami wariantu ogólnego zawierają łącznie $6\cdot0{,}75+1=5{,}5$ razy tyle wag co pojedynczy bazowy MLP, bez uwzględnienia routera. \newline
Oszczędność dotyczy warunkowo wykonywanych obliczeń. Jeżeli udział tokenów skierowanych do ekspertów wynosi $c$, średnia względna szerokość używanego MLP jest równa:
\[
\frac{H_{\mathrm{eff}}}{H}=(1-c)+0{,}75c=1-0{,}25c.
\]
Miara ta przybliża koszt projekcji wyłącznie w zmodyfikowanym MLP. Nie jest redukcją liczby operacji całego modelu ani pomiarem przyspieszenia: router, grupowanie tokenów i wykonywanie małych, nieregularnych partii wprowadzają dodatkowy narzut. \newline
Do odtwarzania rozwiązania przygotowano manifest zawierający identyfikację modelu bazowego, maski, checkpoint routera, kolejność domen i próg pewności. Rozmiary plików oraz sumy kontrolne SHA-256 są weryfikowane przed złożeniem modelu. Wariant ten rekonstruuje ekspertów bezpośrednio z masek i wag bazowych, dzięki czemu nie wymaga zapisywania osobnej pełnej kopii wag każdego eksperta na dysku; kompaktowe kopie nadal powstają w pamięci podczas uruchomienia. Konfigurację masek i routera zamrożono przed końcową ewaluacją. Uzyskane rozwiązanie pozwala zatem zbadać, czy wybór neuronów oparty na domenach SAE zachowuje użyteczne zachowanie modelu przy mniejszej aktywnej szerokości jednej warstwy. Odpowiedź na to pytanie wymaga porównania z modelem bazowym i kontrolnymi sposobami selekcji, przedstawionego w kolejnym rozdziale.
